\section{Sampling and deterministic normalization}
\label{sec:measured-convergence}
\label{measure:section}

We use two elementary ways to retain the sampling measure when a metric
limit is identified. The first uses finite distance arrays and covering
radii; the second uses a common time parametrization. We then give one
deterministic normalization argument for both models.

\subsection{From sampled distances to measured convergence}

A compact metric probability space is a triple $(X,d,\mu)$, where
$(X,d)$ is a nonempty compact metric space and $\mu$ is a Borel probability
measure. We identify triples by measure-preserving isometries. Its
Gromov--Hausdorff--Prokhorov distance from $(Y,e,\nu)$ is the infimum,
over isometric embeddings in a common metric space, of the maximum of
the Hausdorff distance between their images and the Prokhorov distance
between their measures. The latter distance is the infimum of $r>0$
such that $\mu(A)\leq\nu(A^r)+r$ and $\nu(A)\leq\mu(A^r)+r$
for every Borel set $A$, where $A^r$ is its open $r$-neighborhood.
The space of these isometry classes is separable and complete in this
metric \cite[Theorem~6 and Proposition~8(i)]{MiermontTessellations}.

Sampled distance arrays underlie the Gromov-weak framework of
Greven, Pfaffelhuber, and Winter \cite{GrevenPfaffelhuberWinter}.
Controlling the full support requires an additional condition, as in
the lower-mass criterion of Athreya, L\"ohr, and Winter
\cite{AthreyaLohrWinter}. We use the following covering-radius version
and give its proof. The covering assumption rules out points which carry
little mass but remain far from all typical samples.

\begin{lemma}[Sample arrays and covering radii]\label{measure:criterion}
Let $(X_n,d_n,\mu_n)$ and $(X,d,\mu)$ be random compact metric
probability spaces, with $\mu$ almost surely of full support.
Conditional on each space, sample a sequence $(x_{n,i})_{i\geq1}$,
respectively $(x_i)_{i\geq1}$, independently with its probability measure.
Suppose that, for every fixed $k$,
\[
 (d_n(x_{n,i},x_{n,j}))_{i,j\leq k}
 \ \Longrightarrow\ (d(x_i,x_j))_{i,j\leq k},
\]
and that, for every $r>0$,
\begin{equation}\label{measure:covering}
 \lim_{k\to\infty}\limsup_{n\to\infty}
 \mathbb P(R_{n,k}>r)=0,
 \qquad
 R_{n,k}:=\sup_{z\in X_n}\min_{i\leq k}d_n(z,x_{n,i}).
\end{equation}
Then $(X_n,d_n,\mu_n)$ converges in distribution to $(X,d,\mu)$
in Gromov--Hausdorff--Prokhorov topology.
Only the indicated conditional sampling laws are assumed; no joint
coupling of the spaces is needed.
\end{lemma}

\begin{proof}
Write $\widehat\mu_{n,m}=m^{-1}\sum_{i=1}^m\delta_{x_{n,i}}$.
We first prove the uniform empirical approximation
\begin{equation}\label{measure:empirical-uniform}
 \lim_{m\to\infty}\limsup_{n\to\infty}
 \mathbb P\bigl(d_{\mathrm P}^{X_n}(\mu_n,\widehat\mu_{n,m})
                     >\varepsilon\bigr)=0
 \quad(\varepsilon>0).
\end{equation}
Use the next $k$ samples as auxiliary centers
$y_{n,j}=x_{n,m+j}$. They are independent of the first $m$ samples
conditional on the space. Partition $X_n$ into the Borel cells
$C_1,\ldots,C_k$ of the nearest center, breaking ties by the smallest
index. Conditional on the space and the centers, put
\[
 p_j=\mu_n(C_j),\qquad
 \widehat p_j=\frac1m\sum_{i=1}^m\mathbf1_{\{x_{n,i}\in C_j\}},
 \qquad T=\frac12\sum_{j=1}^k|\widehat p_j-p_j|.
\]
The conditional multinomial law and Cauchy--Schwarz give
\begin{equation}\label{measure:cell-counts}
 \mathbb E[T\mid X_n,d_n,\mu_n,(y_{n,j})_{j\leq k}]
 \leq \frac1{2\sqrt m}\sum_{j=1}^k\sqrt{p_j(1-p_j)}
 \leq\frac12\sqrt{\frac{k}{m}}.
\end{equation}
If these centers form an $r$-net, moving every point to its cell's
center moves both $\mu_n$ and $\widehat\mu_{n,m}$ by distance at most
$r$. The two resulting measures on the centers have total variation
distance at most $T$, and hence Prokhorov distance at most $T$.
Indeed their common mass can be coupled at identical centers and the
remaining mass is at most $T$. Thus
\[
 d_{\mathrm P}^{X_n}(\mu_n,\widehat\mu_{n,m})\leq 2r+T.
\]
The auxiliary covering radius has the same law as $R_{n,k}$. Markov's
inequality therefore yields, for $r,t>0$,
\begin{equation}\label{measure:empirical-bound}
 \mathbb P\bigl(d_{\mathrm P}^{X_n}(\mu_n,\widehat\mu_{n,m})
                         >2r+t\bigr)
 \leq \mathbb P(R_{n,k}>r)+\frac1{2t}\sqrt{\frac{k}{m}}.
\end{equation}
Choose $r,t$ with $2r+t<\varepsilon$, then $k$ using
Equation~\eqref{measure:covering}, and finally let $m$ tend to infinity.
This proves Equation~\eqref{measure:empirical-uniform}. The independence
of the auxiliary centers is essential: the cells were fixed before
the sample counts in Equation~\eqref{measure:cell-counts} were drawn.

Let $F_{n,m}$ be the finite metric space consisting of the first $m$
sampled points, equipped with $\widehat\mu_{n,m}$, with repeated points
identified and their masses added. Inside $X_n$ its Hausdorff distance
from $X_n$ is exactly $R_{n,m}$. Consequently
\begin{equation}\label{measure:finite-approximation}
 \lim_{m\to\infty}\limsup_{n\to\infty}
 \mathbb P\bigl(d_{\mathrm{GHP}}((X_n,d_n,\mu_n),F_{n,m})
                                      >\varepsilon\bigr)=0.
\end{equation}

For fixed $m$, the distance-array assumption gives $F_{n,m}\Rightarrow F_m$,
where $F_m$ is the analogous empirical space from $x_1,\ldots,x_m$.
Here is the needed continuity, including possible repeated points.
Two $m$-label distance arrays differing entrywise by at most $a$ give
a correspondence pairing equal labels with distortion at most $a$.
For each $b>a/2$, join paired points in the disjoint union by edges of
length $b$ and take the induced path metric. The distortion inequality
ensures that no new path shortens either original metric. Paired
points are at distance at most $b$, so the Hausdorff distance is at most
$b$; the coupling which gives mass $1/m$ to each label pair makes the
Prokhorov distance at most $b$ as well. Letting $b\downarrow a/2$
proves the claimed continuity. Zero distances simply identify labels.

Almost surely the target samples are dense: every nonempty ball in a
countable base has positive measure and is eventually sampled.
Compactness then implies $R_m:=\sup_{z\in X}\min_{i\leq m}d(z,x_i)\to0$.
The same auxiliary-center argument, now for this single random space,
shows that its empirical Prokhorov distances tend to zero in probability.
Therefore $F_m\to(X,d,\mu)$ in probability in GHP distance.

For completeness, let $f$ be bounded by one and Lipschitz with constant
one for $1\wedge d_{\mathrm{GHP}}$. Insert $F_{n,m}$ and $F_m$ between
the two arguments of $f$. For fixed $m$ the middle expectation
difference tends to zero. The other differences tend to zero after
$m\to\infty$ by Equation~\eqref{measure:finite-approximation} and the
target approximation. Bounded Lipschitz functions determine weak
convergence on the space of compact metric probability spaces. This
proves the result.
\end{proof}

\begin{lemma}[The measure in a common sample completion]
\label{measure:common-completion}
Suppose two random compact metric probability spaces with full-support
measures have sample sequences which are conditionally iid in each
marginal. If their labeled distance arrays agree almost surely, their
unique label-preserving completion isometry preserves their measures.
No conditional independence given both spaces is required.
\end{lemma}
\begin{proof}
In each marginal, almost surely the samples are dense and their empirical
measures converge weakly to the sampling measure. For the latter
assertion, condition on a fixed compact space and apply the strong law
to a countable uniformly dense subset of its continuous functions.
The two probability-one events intersect under any coupling.
The equal arrays define an isometry of the dense sample sets, hence a
unique isometry of the compact completions. It maps every empirical
measure in one space to the corresponding empirical measure in the
other. Passing to weak limits proves the assertion.
\end{proof}

\begin{lemma}[A common parameter and its measure]\label{measure:parameter}
Let $f_n:[0,1]\to X_n$ and $f:[0,1]\to X$ be measurable surjections
onto compact metric spaces, and let their pushforwards of Lebesgue
measure be $\mu_n$ and $\mu$. If
\[
 e_n:=\sup_{s,t\in[0,1]}
 |d_n(f_n(s),f_n(t))-d(f(s),f(t))|\longrightarrow0,
\]
then $d_{\mathrm{GHP}}((X_n,d_n,\mu_n),(X,d,\mu))\leq e_n/2$.
\end{lemma}
\begin{proof}
The correspondence $\{(f_n(t),f(t)):t\in[0,1]\}$ covers both spaces
and has distortion at most $e_n$. Use the common metric constructed in
the proof of Lemma~\ref{measure:criterion}, with bridges of length $b>e_n/2$.
The pushforward of Lebesgue measure under $t\mapsto(f_n(t),f(t))$
is a coupling of the two measures, supported on pairs at distance at
most $b$. Both Hausdorff and Prokhorov distances are at most $b$.
Let $b\downarrow e_n/2$.
\end{proof}

\subsection{Removing a deterministic subsequential multiplier}

\begin{lemma}[A continuous translation-equivariant center]
\label{measure:center}
\label{tree:lem:diameter_center}
For every finite real random variable $Y$, there is a unique $m(Y)\in\mathbb R$
satisfying
\begin{equation}\label{tree:eq:sphere_arctan_center}
                    \mathbb E\arctan(Y-m(Y))=0.
\end{equation}
It depends only on the law of $Y$, satisfies $m(Y+c)=m(Y)+c$ for every
deterministic $c\in\mathbb R$, and is continuous under convergence in
distribution of finite real random variables.
\end{lemma}
\begin{proof}
For fixed $Y$, the function $t\mapsto\mathbb E\arctan(Y-t)$ is
continuous, strictly decreasing, and has limits $\pi/2$ and $-\pi/2$.
Thus its zero exists uniquely. Translation equivariance follows by
substitution. If $Y_j\Rightarrow Y$, expectations converge at every fixed
$t$, since the integrand is bounded and continuous. Their strict opposite
signs at $m(Y)\pm\varepsilon$ trap the zeros, so $m(Y_j)\to m(Y)$.
No moment assumption is used.
\end{proof}

\begin{lemma}[Deterministic normalization of a measured limit]
\label{measure:normalization}
Let $(V_n,d_n,\mu_n)$ be random compact metric probability spaces, let
$s_n>0$ be deterministic with $s_n\to\infty$, and let $(S,D,\mu)$ be a random
compact metric probability space whose diameter $H$ is finite and
strictly positive almost surely. Suppose every subsequence has a
further subsequence on which
\[
                 (V_n,d_n/s_n,\mu_n)\Longrightarrow(S,cD,\mu)
\]
in GHP topology, for some deterministic $c\in(0,\infty)$ which may
depend on that further subsequence. With $m$ as in Lemma~\ref{measure:center},
write $H_n=\operatorname{diam}(V_n,d_n)$ and put
\[
 b_n=\exp\{m(\log(H_n\vee1))-m(\log H)\},\qquad a_n=b_n^{-1}.
\]
Then $a_n>0$, $a_n\to0$, and
\[
                  (V_n,a_nd_n,\mu_n)\Longrightarrow(S,D,\mu)
\]
along the full sequence.
\end{lemma}
\begin{proof}
Take any subsequence and a further subsequence from the hypothesis.
Diameter is continuous for GH, and GHP convergence implies GH
convergence, so $H_n/s_n\Rightarrow cH$. Since $s_n\to\infty$ and
$cH>0$ almost surely,
\[
 \log(H_n\vee1)-\log s_n\Rightarrow\log c+\log H.
\]
Continuity and translation equivariance of $m$ give $b_n/s_n\to c$.
Multiplying the metrics by the positive deterministic ratios
$s_n/b_n\to c^{-1}$ gives the stated normalized GHP limit on this
further subsequence, with the measures unchanged. Positive scaling is
jointly continuous in GHP: scale a common embedding for fixed factors;
for convergent factors, the identity correspondence has distortion
bounded by their difference times the diameter, and the identity
coupling controls the measures in the same embedding.
Every subsequence has such a further subsequence, which proves the
full-sequence assertion. If $b_n$ failed to tend to infinity, a bounded
subsequence of $b_n$ would have a further subsequence with
$b_n/s_n\to c>0$, a contradiction.

All centers here use unconditional diameter laws, so the normalization
is deterministic. The truncation $H_n\vee1$ includes one-vertex maps.
The hypothesis that $c$ is deterministic is indispensable: this
argument does not remove a random multiplier.
\end{proof}
